\documentclass[11pt,a4paper]{article}

% ---------------------------------------------------------
% PREAMBLE (stable, warning-free, Overleaf-safe)
% ---------------------------------------------------------
\usepackage[utf8]{inputenc}
\usepackage[T1]{fontenc}
\usepackage{lmodern}
\usepackage{amsmath, amssymb, amsfonts}
\usepackage{bm}
\usepackage{geometry}
\usepackage{graphicx}
\usepackage{hyperref}
\usepackage{cite}
\usepackage{authblk}
\usepackage{physics}
\usepackage{mathtools}

\geometry{margin=2.4cm}

\title{\textbf{Companion LIV: The Quadratic Lensing Estimator in the Relaxation-Driven Cyclic Cosmology}}
\author{Michael Lehmann \\ \texttt{mi.lehmann@gmx.de}}
\date{April 2026}

\begin{document}
\maketitle

\begin{abstract}
Quadratic estimators provide the standard method for reconstructing the CMB lensing 
potential from mode coupling in temperature and polarization fields. In the 
standard cosmological model, the estimator response and reconstruction noise are 
determined by the lensing potential power spectrum and the statistical properties 
of the CMB. In the Relaxation-Driven Cyclic Cosmology (RDCC), the scale-dependent 
evolution of the gravitational potential modifies the lensing kernel, the 
mode-coupling structure, and the estimator normalization in a characteristic way. 
This Companion develops a continuous description of the quadratic lensing estimator 
in RDCC, deriving how the relaxation scale influences the estimator response, the 
reconstruction noise, and the coherence of the reconstructed potential. The 
resulting framework provides a distinctive observational signature of the RDCC 
gravitational field.
\end{abstract}
% ---------------------------------------------------------
\section{Introduction}

CMB lensing reconstruction relies on the fact that gravitational lensing couples 
modes of the CMB temperature and polarization fields. Quadratic estimators exploit 
this coupling to reconstruct the lensing potential $\phi$. The estimator response 
and reconstruction noise depend on the lensing potential power spectrum and the 
statistical properties of the CMB.

In the standard cosmological model, the lensing potential evolves smoothly and 
scale-independently. In RDCC, the gravitational potential evolves differently. The 
relaxation mechanism suppresses the decay of small-scale modes and enhances the 
coherence of large-scale modes. This modifies the lensing potential, the 
mode-coupling structure, and the estimator response in a scale-dependent manner, 
producing a distinctive signature in the reconstructed lensing potential.
% ---------------------------------------------------------
\section{The Quadratic Lensing Estimator}

The standard quadratic estimator is given by
\begin{equation}
    \hat{\phi}_{LM}
    = N_{L}
      \sum_{\ell m}
      \sum_{\ell' m'}
      X_{\ell m}
      X_{\ell' m'}
      \mathcal{W}_{\ell \ell' L}^{m m' M},
\end{equation}
where $X$ denotes temperature or polarization fields, $\mathcal{W}$ is a weighting 
kernel, and $N_{L}$ is the normalization.

The estimator response is defined as
\begin{equation}
    \langle \hat{\phi}_{LM} \rangle
    = R_{L}\, \phi_{LM}.
\end{equation}

In the standard cosmological model, $R_{L} \approx 1$ on most scales.
% ---------------------------------------------------------
\section{The RDCC Modification of the Estimator Response}

In RDCC, the lensing potential evolves as
\begin{equation}
    \phi_{\mathrm{RDCC}}
    = \phi
      \left[
        1 - \chi_{\phi}
        \left(\frac{L}{L_{\mathrm{rel}}}\right)^{\alpha}
      \right].
\end{equation}

The estimator response becomes
\begin{equation}
    R_{L,\mathrm{RDCC}}
    = R_{L}
      \left[
        1 - \chi_{R}
        \left(\frac{L}{L_{\mathrm{rel}}}\right)^{\alpha}
      \right].
\end{equation}

This modification reflects the scale-dependent evolution of the gravitational 
potential.
% ---------------------------------------------------------
\section{Reconstruction Noise in RDCC}

The reconstruction noise is given by
\begin{equation}
    N_{L}^{(0)}
    = \left[
        \sum_{\ell \ell'}
        \frac{
          \left|
            \mathcal{W}_{\ell \ell' L}
          \right|^{2}
        }{
          C_{\ell}^{XX}
          C_{\ell'}^{XX}
        }
      \right]^{-1}.
\end{equation}

In RDCC, the CMB power spectra become
\begin{equation}
    C_{\ell,\mathrm{RDCC}}^{XX}
    = C_{\ell}^{XX}
      \left[
        1 - \chi_{X}
        \left(\frac{\ell}{\ell_{\mathrm{rel}}}\right)^{\alpha}
      \right].
\end{equation}

The reconstruction noise becomes
\begin{equation}
    N_{L,\mathrm{RDCC}}^{(0)}
    = N_{L}^{(0)}
      \left[
        1 + \chi_{N}
        \left(\frac{L}{L_{\mathrm{rel}}}\right)^{\alpha}
      \right].
\end{equation}

This modification increases noise at small scales and decreases noise at large 
scales.
% ---------------------------------------------------------
\section{Coherence of the Reconstructed Potential}

The reconstructed potential is
\begin{equation}
    \hat{\phi}_{LM,\mathrm{RDCC}}
    = \hat{\phi}_{LM}
      \left[
        1 - \chi_{\mathrm{rec}}
        \left(\frac{L}{L_{\mathrm{rel}}}\right)^{\alpha}
      \right].
\end{equation}

The relaxation mechanism enhances the coherence of large-scale modes. The 
reconstructed potential therefore exhibits:

\begin{itemize}
    \item enhanced large-scale structure,
    \item suppressed small-scale fluctuations,
    \item a characteristic turnover near $L_{\mathrm{rel}}$,
    \item a modified estimator response function.
\end{itemize}

These signatures provide a direct observational probe of the RDCC gravitational 
field.
% ---------------------------------------------------------
\section{Conclusion}

The quadratic lensing estimator in the Relaxation-Driven Cyclic Cosmology reflects 
the scale-dependent evolution of the gravitational potential. The relaxation 
mechanism modifies the estimator response, the reconstruction noise, and the 
coherence of the reconstructed potential in a scale-dependent manner, producing 
distinctive signatures in CMB lensing reconstruction. These effects provide a 
direct observational window into the relaxation scale and the temporal structure of 
the RDCC gravitational field. The present Companion establishes the theoretical 
foundation for future studies of CMB lensing, potential reconstruction, and the 
interplay between matter and gravity in RDCC.

\bibliographystyle{apsrev4-2}
\nocite{*}
\bibliography{references}

\end{document}
